This paper examines whether providing households with digital feedback on their electricity use affects consumption behavior. Using a framed field experiment with high-frequency smart meter data, we find that access to a feedback application leads to a reduction in electricity consumption of around 4.5--4.7\%. The effect emerges relatively quickly after rollout and remains persistent for several months. The adjustment is primarily observed during hours in which households have greater scope to modify their electricity use. At the same time, the response is neither uniform nor stable over time. Point estimates remain negative throughout the observation period, but statistical significance is intermittent rather than continuous: the reduction is significant in most weeks through week 25 and only sporadically thereafter, without a sharp cutoff.

A key insight of the study concerns the role of attention. Engagement with the application declines markedly after the initial rollout and is concentrated among a relatively small group of users. The top 10\% of users account for 76.7\% of all recorded interactions. Households that interact more frequently with the app exhibit stronger and more persistent reductions in consumption, whereas households with little or no engagement show only short-lived changes. These patterns point toward the importance of attention and repeated interaction, although engagement itself is not exogenous and may reflect underlying differences between households.

Taken together, the findings suggest that information provision can influence behavior, but that its effectiveness depends on continued interaction with the information provided. Without such interaction, the impact appears to diminish over time. At the same time, the fact that consumption does not immediately return to baseline levels indicates that responses may not be purely transitory.

Several limitations of the study should be acknowledged. First, the analysis identifies the effect of being offered access to the application, but not the causal effect of actually using it. Since engagement is a post-treatment outcome, differences across users may be driven by unobserved characteristics such as motivation or environmental concern. Second, although the dynamic analysis provides insight into how effects evolve over time, it does not allow us to clearly distinguish between different behavioral mechanisms, such as increased salience, learning, or the formation of habits.

Third, participation required active registration and installation of a smartphone application, which may limit the generalizability of the findings to less engaged or less technologically inclined populations. Finally, while we observe reductions in electricity consumption, we cannot directly assess whether these changes improve welfare. Lower consumption may reflect more informed decision-making, but it may also involve reductions in valued energy services.

Future research could address these limitations by experimentally varying the intensity and timing of feedback, or by combining information provision with other instruments such as pricing incentives or automated technologies. Such approaches may help to better understand how behavioral responses can be sustained over longer periods.

From a policy perspective, the results suggest that information-based interventions can contribute to reducing residential electricity demand. However, their effectiveness appears to depend on maintaining user attention, which may require complementary design features or additional policy instruments.